\section{Introduction}
\label{sec:introduction}
Prior authorization (PA) is a utilization management process requiring providers to obtain insurer approval before rendering certain services. Insurers adjudicate requests by approving, modifying, denying, or requesting documentation; providers may appeal through successive review levels. This creates substantial administrative burden: physicians complete 39 PA requests weekly, spending 13 hours on PA-related tasks \citep{ama2024pa}. Both sides are now deploying artificial intelligence to gain advantage: insurers use AI to scrutinize claims and increase denials, while providers deploy AI to generate appeals and optimize documentation \citep{williams2024hfma}.

\citet{moritz2025lancet} argue that this dynamic resembles a prisoner's dilemma, warning that ``each party's rational decision to adopt AI for self-protection creates collectively harmful outcomes for everyone.'' The incentive structure is clear: if one party deploys AI, the other must respond or face disadvantage. Yet mutual AI deployment may produce worse outcomes than mutual restraint: more administrative burden, more denials, more appeals, with no improvement in patient care. This raises a fundamental question: what happens when both provider and insurer use large language models to interpret ambiguous coverage policy? Does the predicted coordination failure actually emerge, and through what mechanism?

We address this through a controlled simulation study. Our contributions are:
\begin{enumerate}
    \item \textbf{Multi-agent PA simulation framework.}
    We develop an environment where LLM-driven provider and insurer agents negotiate prior authorization through iterative submission, adjudication, and appeal cycles. Agents operate under information asymmetry (clinical guidelines vs.\ coverage policy) with behavior controlled via prompt-based strategy injection (cooperate, defect, or no guidance).
    \item \textbf{3$\times$3 factorial experiment.}
    On a single gray-zone clinical case, we cross three provider strategies with three insurer strategies and identify a \texttt{pending\_info} vs.\ denial mechanism through which gray-zone ambiguity is resolved. Default (no-guidance) insurer agents apply policy literally, denying gray-zone treatments without cooperative interpretation.
    \item \textbf{Game-theoretic payoff analysis.}
    We derive utility functions from simulation outcomes and show the game's classification is parameter-dependent: at observed administrative costs, a 116$\times$ cost asymmetry makes denial nearly costless for insurers, producing an exploitation structure; under richer cost models, the game can transition to the prisoner's dilemma hypothesized.
\end{enumerate}
